\documentclass[12pt]{article}

\usepackage[a4paper,margin=1in]{geometry}
\usepackage{fontspec}

\defaultfontfeatures{Ligatures=TeX,Renderer=HarfBuzz}
\setmainfont{Latin Modern Roman}
\setmonofont{FreeSerif}
\newfontfamily\banglafont[
  Script=Bengali,
  Renderer=HarfBuzz
]{Noto Serif Bengali}
\newfontfamily\banglasans[
  Script=Bengali,
  Renderer=HarfBuzz
]{Noto Sans Bengali}

\newcommand{\bn}[1]{{\normalfont\banglafont #1}}
\newcommand{\bnsf}[1]{{\normalfont\banglasans #1}}

\title{Bangla LaTeX Starter}
\author{}
\date{\today}

\begin{document}

\maketitle

\section*{English}
This folder is configured for LuaLaTeX with Bangla font support.

\section*{Bangla Example}
\bn{এটি একটি বাংলা উদাহরণ। এখানে আপনি সহজে বাংলা লিখতে পারবেন।}

\section*{Mixed Text}
\bn{আপনি একই ডকুমেন্টে বাংলা ব্যবহার করতে পারবেন।}
You can also write English in the same document.

\section*{How to Write Bangla}
Use \texttt{\textbackslash bn\{...\}} for Bangla serif text.

Example:

\begin{verbatim}
\bn{এখানে বাংলা লিখুন}
\end{verbatim}

\bn{উদাহরণ: আমার বাংলা লেখা}

\bnsf{এটি স্যান্স বাংলা ফন্টে লেখা।}

\end{document}
